\documentclass{jsarticle}
\usepackage[dvipdfmx,final]{graphicx}
\usepackage{amsmath}
\usepackage{amssymb}
\begin{document}
\title{Abandust judgement on economy strategiest \\
built with acceptant of movement \\
Money demand on law of judgement}
\author{Masaaki Yamaguchi, and investigate from AYA TAKASHIMA \\
also future from my son being have with \\
quantum level on differential geometries}

\date{}
\maketitle


高島彩さんの研究論文をアカシックレコードと瞑想法で見ての書き出しとしての
論文


   Retry for future to one result with adapt from adulsent for
aquirance of trainning. For must friends of money are treat on
law of judgement, economy need to treat for aquirered of demands,
this demands supply from rest of  one's suppliement.
This balance end without a closed accidennt of out come for
restorance with supplied of demands exceed.
These system are constructed with Actor, Demand, Supply of
being retrayed from Jones manifold around economy of zone
with accesority, verisity of curve in up, down, eternal
rout of stages.

    This system also monument with partial and economy of
similist from law of judgement in universe and money, economy,
human society, estrand of moment in last eternal stage.
  

   Jones manifold equation built with Euler equation and Euler product,
this concernd with concept from Higgs and partial fields,
moreover this concernd of circle and Euler product from
Global integrate and Volume manifold in Sheap manifold of integrate
with group and toplogy different equation.

$$ e^{- \theta} + e^{i \theta} = \sin (i x \log x) = e^{-f} + e^{f} $$
$$ + \cos (i x \log x) = e^{f} - e^{-f} $$

  自動車における道路での速度とそれに対しての加速度を
  平地と上り坂、下り坂における曲がり角での曲率$R_{ij}$が、
  広中平祐先生の庭園理論と同じテーゼで、彩さんが
  経済理論の景気回復と景気刺激剤、景気恐慌が
  どのようにして、自動車運転と自動車の交通規則と
  同じ理論で作用するかを述べているのが、この論文で説明されている。 


  $$ {d \over dx}\left({d \over d r}R_{ij} = \sin
	  \begin{vmatrix} {l_2 \over {2 \pi}} 
- {l_1 \over {2 \pi}}\end{vmatrix}  < 1 \right) = 0 $$


      Curveの度合い $R_{ij}$、潮汐力の差、${d \over d f}F$とは、
   直接的には違うが、間接的には同じであり、rとRの大域的変数
   として、半径rの多様体が曲率的には$R_{ij}$に大域的微分で作用をしている。
   遠隔的に作用している。


      グリーシャ先生からリッチ・フロー方程式が、大域的微分になるのを
   教わった。広中平祐先生の4重帰納法で、ガンマ関数における大域的
   微分多様体が、オイラーの定数の多様体積分に使えるのがわかった。
   これがヒルベルト空間としてのサーストン多様体が、
   代数多様体としての微分幾何として、ゼータ関数が素数分布として
   展開されるのがわかった。

  $$ R_{ij} = |R_2 - R_1| < 1, |R_2 - R_1| = 0 $$
  $$ R_{ij} = {R_2 \over R_1} < 1, {R_2 \over R_1} = 1 $$
  $$ {d \over dt}g_{ij} = - 2 R_{ij}, {d \over d f} F(x) = 
  \int \Gamma(\gamma)^{'}dx_m = \int \Gamma dx_m + {d \over d \gamma}
  \Gamma \le e^{f} - e^{-f} \le e^{-f} + e^{f} $$

\includegraphics[width=6cm,clip]{curver1beforer2}
\includegraphics[width=6cm,clip]{curver1r2}
\includegraphics[width=6cm,clip]{curver1tor2}
\includegraphics[width=6cm,clip]{curver2tor1}
\includegraphics[width=6cm,clip]{curvr1r2}

$$ R_2 < R_1 \text{の場合は、}$$
右からR1に接近して、R2へ行くが、${R_2 \over R_1} < 1 $
より、曲率$R_{ij} < 1$であるから、平坦の道路では、気をつけるべし、
曲率$R_{ij} = 1$では、CURVEには気をつける。
曲率$R_{ij} > 1$では、$R_{i > j}$であり、iの方で気をつける。
これを経済の日経平均相場で、景気の上がりでは、curveがどう分類されるかで、
気をつけるべき時期がわかり、景気の下がりでは、同じくcurveでどう分類するかを
きめる。平坦のときが一番気をつける、油断できない景気の税と立法での
判断が、自動車の走行分岐での心理作用が、景気判断のどの道路状況で
その場の気をつける心理と同じ作用をする。
経済の需要が供給を下回っている状況では、$\text{demand} \le \text{supply}$
では、需要のために、物価を下げるが、この状況では、$ R_1 < R_2, $
であるので、株は$R_2$の方へ曲率がまがり、このときに、需要の
企業の生産が減ってきて、このときに、金融機関が気をつけないと、
企業が破綻する可能性がある。このシチュエーションが世界恐慌であった。
このように、曲率の場合分けをしていないと、間違った判断で、
曲率に気をつけないから、交通事故が起こるということである。
広中平祐先生の庭園理論と同じ考えのモチーフを高島彩さんは、
研究論文で提出している。

   高島彩さんの成蹊大学での研究論文を参照させてもらえて、
書き出した経済の需要と供給に対しての物価上昇と経済沈滞
においてのインフレとデフレにメスを入れる税と景気対策
のタイミング、このタイミングがJones多項式の周期の終わり
にバランスを崩すと、景気の極率としてのCURVEの潮汐力での
差分率の曲率が１以下になるときの、心理作用を考えるべき
という、経済消費における物価調整の作用をかけるメカニズム
についての論文 　

参考文献: 父、母、彩さん、ナッシュ先生、益川先生、ワインバーグ先生、まどかさん


\end{document}
